\documentclass[11pt,a4paper]{article}
\usepackage[margin=2.15cm]{geometry}
\usepackage{booktabs}
\usepackage{tabularx}
\usepackage{array}
\usepackage{xcolor}
\usepackage{hyperref}
\usepackage{enumitem}
\usepackage{titlesec}
\hypersetup{colorlinks=true,linkcolor=blue,urlcolor=blue}
\definecolor{evoblue}{HTML}{0071E3}
\definecolor{evogreen}{HTML}{248A3D}
\definecolor{evoamber}{HTML}{C66A00}
\definecolor{evored}{HTML}{D70015}
\newcolumntype{Y}{>{\raggedright\arraybackslash}X}
\titleformat{\section}{\large\bfseries\color{evoblue}}{\thesection}{0.55em}{}
\titleformat{\subsection}{\normalsize\bfseries}{\thesubsection}{0.55em}{}
\setlist[itemize]{leftmargin=1.2em,itemsep=0.25em}

\title{\textbf{Industrial Predictive Quality MVP}\\Plain-English Sales Briefing for Audi and Industrial Manufacturers}
\author{EVOCON Solutions}
\date{\today}

\begin{document}
\maketitle

\begin{center}
\fbox{\begin{minipage}{0.92\linewidth}
\textbf{One sentence:} EVOCON's Industrial Predictive Quality MVP helps factories find visual defects, machine-risk patterns and process deviations earlier by combining images, sensors and process data into a ranked inspection queue.
\end{minipage}}
\end{center}

\section{Purpose of This Document}
This document is written for a non-technical salesperson or executive who needs to explain the MVP in a factory meeting. It avoids academic language and focuses on what the customer cares about:
\begin{itemize}
  \item fewer late defects;
  \item fewer unexpected machine problems;
  \item less noise from alarms;
  \item clearer inspection priorities;
  \item a low-risk pilot that can prove or disprove value in 4--6 weeks.
\end{itemize}

\section{What Problem Are We Solving?}
Modern factories already collect large amounts of data: camera images, sensor traces, PLC/SCADA values, torque curves, test logs, recipes, setpoints and maintenance events. The problem is that these signals are usually reviewed separately.

The MVP turns these separate signals into one practical question:
\begin{quote}
\emph{Which part, cycle, machine, batch or line should the team review first?}
\end{quote}

For a premium automotive manufacturer such as Audi, the relevant pain points are surface quality, assembly quality, machine drift, supplier quality, line stoppages and late detection of anomalies.

\section{Explain It Like I Am Five}
Imagine an experienced quality engineer who sees thousands of good parts and machine cycles. Over time, the engineer develops a feeling for what normal production looks like.

The MVP does something similar, but with data:
\begin{enumerate}
  \item it looks at images and sensor traces from normal production;
  \item it learns what good production usually looks like;
  \item it watches new production data;
  \item it highlights what changed;
  \item it ranks what should be checked first.
\end{enumerate}

It does not replace engineers. It gives engineers a better review queue.

\section{The Product in Plain English}
\begin{tabularx}{\linewidth}{@{}p{0.25\linewidth}Y@{}}
\toprule
\textbf{Module} & \textbf{What it does} \\
\midrule
Visual Quality & Looks at part or product images, marks suspicious areas with heatmaps and ranks images by review priority. \\
Machine Risk & Uses vibration, current, temperature, torque, SCADA or test logs to score whether a machine/cycle is becoming risky. \\
Process Deviation & Compares how a process should evolve with how it actually evolves, then marks the first point where it starts drifting. \\
Hierarchical Ranking & Aggregates local evidence from patch to image, part, cycle, batch and line, so managers see useful priorities instead of raw signals. \\
\bottomrule
\end{tabularx}

\section{Why This Fits Audi}
Audi and Volkswagen Group plants have complex production systems with high quality expectations. The MVP is relevant wherever the plant has:
\begin{itemize}
  \item image-based inspection of parts, surfaces, welds, paint, connectors or assemblies;
  \item torque, vibration, current, temperature or test-log traces;
  \item station-level or line-level quality events;
  \item process settings such as recipe, speed, force, pressure, temperature, material, tool or robot program;
  \item a need to reduce manual review burden or identify the first warning signs of drift.
\end{itemize}

\section{Suggested First Audi Use Cases}
\begin{tabularx}{\linewidth}{@{}p{0.28\linewidth}Y Y@{}}
\toprule
\textbf{Use case} & \textbf{Data needed} & \textbf{Output} \\
\midrule
Body / paint / surface inspection & Images from good and questionable parts, station ID, inspection results & Heatmap and part review ranking \\
Assembly station risk & Torque curves, press-fit traces, robot/PLC events, station ID & Top risky cycles and deviation timeline \\
Tool or machine drift & Vibration/current/temperature, maintenance events, process context & Machine risk score and early-warning ranking \\
Supplier quality & Incoming-part images, batch/supplier metadata, inspection labels & Batch and supplier review prioritisation \\
End-of-line test anomaly & Test logs, final pass/fail labels, process context & Test anomaly score and root-candidate ranking \\
\bottomrule
\end{tabularx}

\section{What Results Can We Honestly Show?}
Current internal MVP results support selling a pilot, not a production guarantee.

\begin{tabularx}{\linewidth}{@{}p{0.30\linewidth}p{0.22\linewidth}Y@{}}
\toprule
\textbf{Area} & \textbf{Current evidence} & \textbf{Commercial meaning} \\
\midrule
Visual anomaly detection & DINOv2/ResNet + PatchCore/PaDiM achieved near-perfect quick-run scores on MVTec bottle. & Visual inspection is the strongest first product module to pilot. \\
Sensor failure-soon risk & Engineered sensor features are strong; future learned signal is promising but not a replacement. & Start with reliable physical indicators, then add learned signals if they improve validation. \\
No-cycle sensor evidence & Sensor and latent signals help when lifecycle/cycle proxies are removed. & Useful when simple counters do not explain risk. \\
Hardness/material generalization & Sensor features were strong in held-out hardness/material-style splits. & Physical signals can matter when process regime changes. \\
\bottomrule
\end{tabularx}

\section{What Not To Claim}
Do not say:
\begin{itemize}
  \item ``This is production-ready tomorrow.''
  \item ``This replaces quality engineers.''
  \item ``The model guarantees zero defects.''
  \item ``The research model always beats engineered features.''
  \item ``We have already deployed this at Audi.''
\end{itemize}

Say instead:
\begin{quote}
``We have a pilot-ready MVP. We can take one process, compare against strong baselines and show in 4--6 weeks whether images, sensors and process data can detect defects or drift earlier than current checks.''
\end{quote}

\section{The Recommended Meeting Script}
\subsection*{Opening}
``EVOCON already works in industrial automation and project support. This MVP adds an AI layer for predictive quality: it combines camera inspection, sensor trends and process context to produce a ranked review queue. We are not asking you to replace your systems. We propose a 4--6 week pilot on one line to measure whether we can catch defects or machine drift earlier.''

\subsection*{First Question}
``Where do you currently detect quality or machine issues too late?''

\subsection*{If They Mention Visual Quality}
``Then the first pilot should be visual: part images in, heatmaps and inspection ranking out. We compare against simple rules and strong visual baselines.''

\subsection*{If They Mention Machine Downtime}
``Then the first pilot should use sensor/process data: vibration, current, temperature, torque, SCADA or test logs. We rank cycles or machines by risk and measure lead time and false alarms.''

\subsection*{If They Mention Too Many Alarms}
``Then the first pilot should focus on prioritisation: the top 5--10\% alerts that deserve human review, with evidence for why they were selected.''

\section{What Data To Ask For}
\begin{tabularx}{\linewidth}{@{}p{0.28\linewidth}Y@{}}
\toprule
\textbf{Minimum} & 100--1,000 historical images or sensor windows from one process, with timestamps and any known quality/failure labels. \\
\midrule
\textbf{Better} & Images + sensor traces + process settings + station IDs + inspection results + maintenance/quality events. \\
\midrule
\textbf{Best} & A full sequence: part/cycle ID, timestamp, station, line, recipe/setpoint, image, sensor trace, inspection result and downstream event. \\
\bottomrule
\end{tabularx}

\section{4--6 Week Pilot Plan}
\begin{tabularx}{\linewidth}{@{}p{0.18\linewidth}Y@{}}
\toprule
\textbf{Week} & \textbf{Work} \\
\midrule
1 & Data audit, target definition and baseline agreement. \\
2 & Build simple baselines and visual/sensor foundation models. \\
3--4 & Train/evaluate risk ranking, heatmaps, deviation timelines and calibration. \\
5 & Review with plant experts; collect false positives, missed cases and practical constraints. \\
6 & Deliver executive report, local demo, metrics, go/no-go recommendation and integration roadmap. \\
\bottomrule
\end{tabularx}

\section{What EVOCON Should Send Before the Meeting}
\begin{itemize}
  \item The ELI5 HTML/PDF deck: \texttt{audi\_industrial\_predictive\_quality\_eli5\_en.html/pdf}.
  \item The algorithm animation: \texttt{assets/industrial\_predictive\_quality\_algorithm\_demo.html}.
  \item A short email, not the full technical paper.
  \item Offer a 15-minute qualification call before asking for data.
\end{itemize}

\section{Short Email Template}
\begin{quote}
Subject: 4--6 week pilot for earlier quality and machine-risk detection

Hello [Name],

EVOCON Solutions works in industrial automation and project support. We are preparing a pilot-ready Industrial Predictive Quality MVP for automotive manufacturing: it combines camera inspection, sensor trends and process context to produce heatmaps, risk scores and a ranked review queue.

The proposal is deliberately narrow: choose one line/process, use historical images or sensor/test data, compare against strong baselines and deliver a go/no-go report in 4--6 weeks.

Would it make sense to have a 15-minute call to identify one quality or machine-risk case where late detection is costly?

Best regards,
[Sender]
\end{quote}

\section{Final Positioning}
The strongest commercial message is:
\begin{quote}
\textbf{We turn existing factory data into earlier warnings and clearer actions.}
\end{quote}

The safest commercial ask is:
\begin{quote}
\textbf{Let us validate one process in a 4--6 week pilot.}
\end{quote}

\end{document}
